\begingroup
\small\setlength{\tabcolsep}{4pt}
\setlength{\TableWidth}{\dimexpr\linewidth-24pt\relax}
\begin{longtable}{@{}P{0.200\TableWidth}P{0.270\TableWidth}P{0.260\TableWidth}P{0.270\TableWidth}@{}}
\caption{Différenciateurs observés dans le benchmark}\label{tab:comparaison}\\
\toprule
\textbf{Critère} & \textbf{Manhattan Active} & \textbf{SAP EWM} & \textbf{Oracle WMS} \\
\midrule\endfirsthead
\multicolumn{4}{l}{\small\itshape Tableau \thetable{} (suite)}\\
\toprule
\textbf{Critère} & \textbf{Manhattan Active} & \textbf{SAP EWM} & \textbf{Oracle WMS} \\
\midrule\endhead
\midrule\multicolumn{4}{r}{\small\itshape Suite à la page suivante}\\\endfoot
\bottomrule\endlastfoot
Idée directrice & Unifier l’exécution & Structurer le contexte métier & Orchestrer le travail \\[3pt]
Point fort le plus visible & Relier stock, tâches, opérateurs, équipements et IA agents & Décrire finement produits, HU, lots, séries, ressources, livraisons & Transformer commande \(\rightarrow\) wave \(\rightarrow\) allocation \(\rightarrow\) tâches \(\rightarrow\) exécution \\[3pt]
Rapport à l’intelligence & IA agents intégrés au plus près de l’action & Intelligence adossée à un modèle métier riche et à l’écosystème SAP & Assistants et recommandations intégrés à l’inventaire, au tasking et aux exceptions \\[3pt]
Apport principal pour ce stage & Boucle courte décision \(\rightarrow\) action & Richesse des données de référence et contraintes produit & Centralité du moteur de tâches et des priorités \\[3pt]
\end{longtable}
\endgroup
